\documentclass[a4paper,12pt]{article}
\usepackage{amsmath}
\usepackage{amssymb}
\usepackage{geometry}
\geometry{margin=1in}

\begin{document}

\title{\textbf{Part - 2: Applying Recursive Spectral Modularity}}
\date{}

\maketitle

\noindent
Let's check if we can further split the communities formed by applying the spectral bipartition method again on each community.

\medskip

We will also compute the \textbf{Degree centrality}, \textbf{Betweenness centrality}, \textbf{Closeness centrality} and \textbf{Clustering coefficient}, which are given by:

\section*{Centrality and Clustering Measures}

\begin{enumerate}
    \item \textbf{Degree centrality:}
    \[
    C_D(i) = \frac{k_i}{n - 1}
    \]
    \texttt{networkx.degree\_centrality(G)}
    
    Measures the fraction of nodes that are directly connected to $i$. Nodes with high degree centrality are immediate hubs with many direct links.

    \item \textbf{Betweenness centrality:}
    \[
    C_B(i) = \sum_{s \neq i \neq t} \frac{\sigma_{st}(i)}{\sigma_{st}}
    \]
    \texttt{networkx.betweenness\_centrality(G)}
    
    Counts how often node $i$ lies on shortest paths between other nodes. Nodes with high betweenness act as bridges or brokers, controlling information flow between different parts of the network.

    \item \textbf{Closeness centrality:}
    \[
    C_C(i) = \frac{n - 1}{\sum_{j \neq i} d(i, j)}
    \]
    \texttt{networkx.closeness\_centrality(G)}
    
    Inverse of the average shortest-path distance from $i$ to all other nodes. A high closeness score indicates that a node can reach all others quickly, making it a good ``spreader'' in the network.

    \item \textbf{Clustering coefficient:}
    \[
    C(i) = \frac{2T(i)}{k_i (k_i - 1)}
    \]
    \texttt{networkx.clustering(G)}
    
    Fraction of a node's neighbors that are also connected to each other. High clustering indicates that $i$ is embedded in a tightly-knit group (a local community or clique).

\end{enumerate}

\end{document}
